% Numerical Setup for SN_A / SN_B Double-Explosion Modelling (concise version)
% SN2024pba Project
%
% Compile: module load texlive/2021
%          pdflatex numerical_setup && bibtex numerical_setup && pdflatex numerical_setup && pdflatex numerical_setup

\documentclass[11pt, a4paper]{article}

\usepackage[round, sort, compress]{natbib}
\usepackage{amsmath, amssymb}
\usepackage{graphicx}
\usepackage{booktabs}
\usepackage{hyperref}
\usepackage[T1]{fontenc}
\usepackage{microtype}

\newcommand{\Msun}{\ensuremath{M_\odot}}
\newcommand{\foe}{\ensuremath{10^{51}\;\mathrm{erg}}}
\newcommand{\kms}{\ensuremath{\mathrm{km\,s^{-1}}}}
\newcommand{\Ni}{\ensuremath{^{56}\mathrm{Ni}}}
\newcommand{\EKO}{\ensuremath{E_{\mathrm{KO}}}}
\newcommand{\Ncnd}{\ensuremath{N_{\mathrm{cnd}}}}

\begin{document}

\title{Numerical Setup}
\author{S. Brennan\\
\small Max Planck Institute for Astrophysics, Garching, Germany}
\date{}

\maketitle

\section*{Numerical Setup}

We model the double-explosion scenario with a two-stage pipeline coupling the
multi-group radiation hydrodynamics code STELLA \citep{Blinnikov1998,
Blinnikov2006} with the non-LTE radiative transfer code CMFGEN
\citep{Hillier1998, Dessart2015}.  In the first stage (SN\_A), a Type~IIb
progenitor --- the solar-metallicity $12\,M_\odot$ model \texttt{s12.0} of
\citet{WoosleyHeger2007}, stripped to a $\sim$3\,$\Msun$ helium core and
matched to model~12C of \citet{Ergon2022} for SN\,2011dh --- is exploded with a
thermal bomb of $0.9 \times 10^{51}$\,erg deposited over 10\,s in the inner
0.3\,$\Msun$, yielding a final kinetic energy of $\sim$0.7\,foe.  A thin
hydrogen envelope of 0.03--0.10\,$\Msun$ ($X = 0.7$) is added to the surface,
and a low-density ($\dot{M} = 10^{-6}\,\Msun\,\mathrm{yr^{-1}}$) sigmoid-blended
circumstellar wind extends to $\sim$$10^{17}$\,cm for numerical stability.  A
\Ni\ mass of 0.075\,$\Msun$ \citep{Ergon2022} is distributed over the inner
0.3\,$\Msun$.  STELLA evolves the explosion with 500 Lagrangian zones using a
GEAR integrator ($\epsilon = 3 \times 10^{-3}$), with $\EKO = 0$ to preserve the
pre-SN velocity structure, $\mathtt{FitTau} = 5$ for frequency-dependent
opacity, and a minimum of $N_{\mathrm{zon}}/10$ diffusion zones to prevent
late-time $\Ncnd$ collapse.

In the second stage (SN\_B), the SN\_A ejecta is homologously evolved to
$\sim$400\,d and serves as the CSM.  A H/He shell (1.0\,$\Msun$, $X = 0.7$,
$Y = 0.3$) with a power-law density profile ($\rho \propto r^{-8}$) following
\citet{Dessart2016} is placed at the centre and launched at 5000\,$\kms$ into
the CSM.  STELLA evolves the shell--CSM interaction for 45\,d with
$\EKO = 0$ (preserving both shell and CSM velocities), enhanced artificial
viscosity ($A_Q = 10$) for the strong velocity reversal, and no \Ni\ (the
luminosity is entirely interaction-powered).

At selected epochs (30\,d post-SBO for SN\_A; 20--40\,d post-launch for SN\_B),
snapshots are extracted from the STELLA output and remapped onto CMFGEN depth
grids ($N_D \sim 150$) using an adaptive algorithm that resolves ionisation
fronts and velocity reversals.  CMFGEN computes non-LTE level populations and
emergent spectra with \texttt{VEL\_LAW\,=\,12} (reading the non-homologous
velocity directly), a tridiagonal depth-coupling solver for the non-monotonic
velocity, Ng acceleration disabled to avoid trace-species instability at
ionisation fronts, and temperature fixed to the hydrodynamic values.  The
inner boundary uses a diffusion condition seeded by the STELLA luminosity.
The complete pipeline --- progenitor conversion, STELLA evolution, snapshot
extraction, CMFGEN spectral synthesis, and dependency-chained job submission
--- is orchestrated by a single workflow manager on the MPCDF Viper HPC
cluster.

\bibliographystyle{aasjournal}
\bibliography{references}

\end{document}
